\chapter{Метафизика}
\label {metaphysics}

«Ваши рассуждения отличаются чудесной законченностью. Если не принимать во внимание несоответствие с действительностью, то это грандиозное достижение мысли» \cite{arnold-cm-r}


Никто из современников Планка революции в физике не предвидел, более того, считалось, что эта наука уже выполнила все стоящие перед ней задачи. Это отметил в юбилейном докладе Королевскому обществу 27 апреля 1900 года патриарх британской физики лорд Кельвин. И только два облачка, по его словам, омрачали ясный научный небосклон. Первое облачко — это вопрос, как может Земля без трения и потери энергии двигаться через упругую среду, каковой является светоносный эфир? А второе облачко — это непреодолимые противоречия теории и опыта в вопросе об излучении «абсолютно чёрного тела».

Примерно то же говорил в 1874 году юному студенту Максу Планку профессор Филипп фон Жолли, руководивший отделением физики Мюнхенского университета \cite{maslov-r}, когда Макс обратился к нему с просьбой записать его в число слушателей лекций по теоретической физике. Маститый учёный попытался отговорить юношу: «Молодой человек! Зачем вы хотите испортить себе жизнь, ведь физика как наука в основном завершена. Осталось прояснить несколько несущественных неясных мест. Стоит ли браться за такое бесперспективное дело?!»1. К счастью для физики, студент оказался настойчивым и находчивым: он ответил, что не собирается открывать ничего нового, а хочет только изучить то, что уже известно.

